\documentclass{article}
\usepackage{graphicx} % Required for inserting images

\title{Buenas Prácticas}
\author{Jennier Melissa Giorgi Ortiz}
\date{September 2026}

\begin{document}

\maketitle
\textbf{Instalación de GitHub en Windows\}
\textbf{Primero se debe instalar GitHub en el computador (si aún no se tiene instalado) para poder acceder a la línea de comando de Git Push.}
}1. Se debe acceder al siguiente enlace: https://product.hubspot.com/blog/git-and-github-tutorial-for-beginners.
2. Al entrar en la página, se desliza hacia abajo y en la palabra "Here" que esta en el "Step 0" se da click para acceder a la página de descarga (Figura 1).\textbf{
\begin{figure}
    \centering
    \includegraphics[width=0.5\linewidth]{image.png}
    \caption{Página para instalación de Git en el computador}
    \label{fig:placeholder}
\end{figure}
3. }Al entrar a la página "Git --distributed-even-if-your-workflow-isnt" se debe deslizar hacia abajo hasta encontrar "Installing on Windows" y en esa parte se da click en el primer enlace que aparece en el párrafo: https://git-scm.com/download/win (Figura 2).\textbf{
\begin{figure}
    \centering
    \includegraphics[width=0.5\linewidth]{image1.png}
    \caption{Enlace para ingresar a la página en donde se podrá instalar Git para Windows}
    \label{fig:placeholder}
\end{figure}
4. }De ahí será redirigido a otra página para iniciar la instalación de GitHub en Windows. Apenas abra ese sitio Web, se debe hacer click en la letra azul que dice: "Click here to download the latest (2.55.0(5)) x64 version of Git for Windows" y automáticamente se comenzará a descargar el archivo .exe para realizar la instalación (Figura 3).\textbf{
\begin{figure}
    \centering
    \includegraphics[width=0.5\linewidth]{image2.png}
    \caption{Enlace para descargar la extensión de instalación de Git}
    \label{fig:placeholder}
\end{figure}
5. }Cuando el .exe se termine de descargar, se dirige al lugar donde guardo el archivo y se le da doble click para iniciar la instalación. A todos los pasos de la instalación se le da "NEXT" hasta que llegue a "INSTALL" y se espera a que llegue al 100% y se instale\textbf{
\textbf{6.} Una vez se termine de instalar, en el computador quedará instalada la aplicación "Git Bash", desde esa extensión se comienza a trabajar. Se puede acceder a ella desde la carpeta en donde quedó guardada la instalación o desde el buscador del escritorio (Figura 4).
\begin{figure}
    \centering
    \includegraphics[width=0.5\linewidth]{image3.png}
    \caption{Interfaz gráfica de Git Bash una vez esta instalado en el computador}
    \label{fig:placeholder}
\end{figure}

\textbf{Apertura de la cuenta en GitHub}
\textbf{1. }Se ingresa a la página de GitHub: github.com (Figura 5)
\begin{figure}
    \centering
    \includegraphics[width=0.5\linewidth]{image4.png}
    \caption{Página para abrir la cuenta en Git Hub}
    \label{fig:placeholder}
\end{figure}

\textbf{2. }En la parte superior derecha de la página web, se da click en "Sign in". La cuenta se puede crear con las cuentas de Google, Apple o con una cuenta aparte en outlook o hotmail. Al dar click en alguna de esas opciones, será redirigido a crear el nombre de usuario y seleccionar el país. Luego se da click en "Create account"
\textbf{3. }Una vez se da click en "Create account" será dirigido a esta página (Figura 6):
\begin{figure}
    \centering
    \includegraphics[width=0.5\linewidth]{image5.png}
    \caption{Página de inicio en Git Hub una vez la cuenta ha sido creada}
    \label{fig:placeholder}
\end{figure}
\textbf{4. }En la esquina superior izquierda se da click en la pestaña verde que dice: "New repository".
\textbf{5. }De ahí será dirigido a la página para crear el repositorio. Ahí, en el espacio que dice "Repository name" se le asigna el nombre al repositorio y lo demás se deja como está por defecto. Al final se da click en "Create repository" (Figura 7).
\begin{figure}
    \centering
    \includegraphics[width=0.5\linewidth]{image6.png}
    \caption{Creación del repositorio en GitHub}
    \label{fig:placeholder}
\end{figure}
\textbf{6. }Ya una vez creado, quedaremos en esta página (Figura 8). Vamos a tener en cuenta el código que aparece ahí en la página para más adelante.
\begin{figure}
    \centering
    \includegraphics[width=0.5\linewidth]{image7.png}
    \caption{Repositorio creado}
    \label{fig:placeholder}
\end{figure}

\textbf{Pasos para clonar la carpeta del profesor Daniel Miranda a nuestro computador}
\textbf{1. }En nuestra cuenta de Git Hub, en la esquina superior derecha hay un espacio que tiene una lupa y dice "Type / To search", ahí se escribe solamente "Dmirandae", no debe decir más. Y eso nos va a dirigir a la siguiente página (Figura 9)
\begin{figure}
    \centering
    \includegraphics[width=0.5\linewidth]{image8.png}
    \caption{Para acceder a la carpeta del profesor Daniel Miranda}
    \label{fig:placeholder}
\end{figure}
\textbf{2. }De ahí, en la barra lateral/izquierda damos click en la pestaña "Users", la cual nos dirigirá a las cuentas que aparecen con ese nombre de usuario (Figura 10). Damos click en la primera que aparece "Daniel R. Miranda-Esquivel".
\begin{figure}
    \centering
    \includegraphics[width=0.5\linewidth]{image9.png}
    \caption{Cuenta del profesor Daniel Miranda}
    \label{fig:placeholder}
\end{figure}
\textbf{3. }Ya dentro del perfil del profesor Miranda, damos click en la carpeta de letras azules que dice: "Dmirandae/nix-BioUIS" e ingresamos a ella (Figura 11).
\begin{figure}
    \centering
    \includegraphics[width=0.5\linewidth]{image10.png}
    \caption{Ingreso a la carpeta Dmirandae/nix-BioUIS}
    \label{fig:placeholder}
\end{figure}
\textbf{4. }Dentro de la carpeta, nos dirigimos a la pestaña de color verde que dice "Code", damos click en ella y copiamos el enlace HTTPS que aparece ahí (Figura 12).
\begin{figure}
    \centering
    \includegraphics[width=0.5\linewidth]{image11.png}
    \caption{Copiar el enlace HTTPS de la carpeta para clonarla desde Git Push}
    \label{fig:placeholder}
\end{figure}

\textbf{Proceso de clonación de la carpeta desde la interfaz Git Push previamente instalada}
\textbf{1. }Ingresamos a Git Bash (Figura 13).
\begin{figure}
    \centering
    \includegraphics[width=0.5\linewidth]{image12.png}
    \caption{Interfaz de Git Bash}
    \label{fig:placeholder}
\end{figure}
\textbf{2. }En la primera línea de comando escribimos "dir" para ver los nombres exactos de los directorios de nuestro computador y elegir en cual vamos a clonar la carpeta de la cuenta del profesor (Figura 14).
\begin{figure}
    \centering
    \includegraphics[width=0.5\linewidth]{image13.png}
    \caption{Directorios del computador}
    \label{fig:placeholder}
\end{figure}
\textbf{3. }Vamos a clonar la carpeta del profesor en la carpeta de "Documents" entonces la siguiente línea de comando sera: "cd Documents" para estar dentro del entorno de la carpeta "Documents" (Figura 15).
\begin{figure}
    \centering
    \includegraphics[width=0.5\linewidth]{image14.png}
    \caption{Enter Caption}
    \label{fig:placeholder}
\end{figure}
\textbf{4. }Una vez dentro del entorno de la carpeta, escribimos la siguiente línea de comando: "git clone https://github.com/Dmirandae/nix-BioUIS.git" y damos enter. Eso nos clonará la carpeta el profesor en nuestra carpeta de "Documents" (Figura 16).
 \begin{figure}
     \centering
     \includegraphics[width=0.5\linewidth]{image16.png}
     \caption{Resultado que debe salir cuando se clona la carpeta}
     \label{fig:placeholder}
 \end{figure}
\textbf{5. }Ahora volvemos al código que nos salió al principio cuando recién creamos nuestra cuenta en GitHub (Figura 8).
\textbf{6. }Ejecutamos tal cual el código en ese orden dentro de la interfaz de Git Bash y por cada línea de comando vamos oprimiendo ENTER (Figura 17).
\begin{figure}
    \centering
    \includegraphics[width=0.5\linewidth]{image17.png}
    \caption{Líneas de comando a seguir}
    \label{fig:placeholder}
\end{figure}
\textbf{7. }Al llegar a la sexta línea de comando (que es un enlace web con el nombre de nuestra cuenta) se solicitará autenticarlo en la cuenta directa de github, entonces automáticamente se abrirá una pestaña en donde se debe dar click en el botón azul que dice "Autenticar desde el navegador" y ahí se redireccionará a la cuenta desde la página web, enviarán un código al correo y al ingresar ese código en la pestaña donde es solicitado, la cuenta quedará verificada y así debe salir en el Git Bash automáticamente sin hacer nada ahí directamente. El programa solo lo hace (Figura 18).
\begin{figure}
    \centering
    \includegraphics[width=0.5\linewidth]{image18.png}
    \caption{Confirmación de autenticación de la cuenta en Git Bash}
    \label{fig:placeholder}
\end{figure}
Y ya con estos pasos quedará instalado Git Push en el computador con sistema operativo Windows. Quedará abierta la cuenta en Git Hub y quedara clonada la carpeta del profesor en la carpeta de nuestro computador en donde decidimos clonarla desde el principio cuando ingresamos en ella con el comando "cd" (Figura 15).







\end{document}
